\documentclass[12pt]{article}
\usepackage[utf8]{inputenc}
\usepackage[margin=1in]{geometry}
\usepackage{amsmath}
\usepackage{amssymb}
\usepackage{natbib}
\usepackage{hyperref}
\usepackage{booktabs}
\usepackage{graphicx}
\usepackage{setspace}
\onehalfspacing

\title{Federal Land Sales as a Source of Treasury Revenue:\\
Policy, Parties, Railroads, and Homesteads, 1790--1900}

\author{Prepared for Thomas J. Sargent and George J. Hall}

\date{\today}

\begin{document}

\maketitle

\begin{abstract}
This paper traces the history of U.S.\ federal land sales as a source of Treasury revenue from the founding era through the end of the nineteenth century.  After Hamilton's funding system consolidated federal fiscal power and ended states' land-for-debt arbitrage schemes, the federal government developed a systematic apparatus for surveying, pricing, and selling the vast public domain.  We document the evolving legislative framework---from the Land Ordinance of 1785 through the Homestead Act of 1862---the partisan disputes over land policy among Federalists, Jeffersonian Republicans, Jacksonian Democrats, Whigs, and the Republican Party, the massive railroad land grants of the 1850s--1870s, and the shift from revenue-maximizing land sales to free distribution under the Homestead Acts.  We provide data on land prices, acreages, revenue magnitudes, and the fiscal consequences of these policy choices.
\end{abstract}

\section{Introduction}

Hamilton's assumption and funding program of 1790--91 resolved the conflict between land speculators and public creditors decisively in favor of creditors, as documented in the companion paper on land speculation and Hamilton's funding system.\footnote{See the companion paper, ``Land Speculation, Public Debt, and Hamilton's Funding System: The Political Economy of Competing Interests, 1783--1795.''}  But that resolution raised an immediate question: what would become of the vast western public domain now that it could no longer serve as a vehicle for retiring depreciated state and Continental debt?

The answer---developed over the next eight decades---made federal land sales one of the two pillars of federal revenue (alongside tariffs) during the long period before the income tax.  The public domain was immense.  Between the original state cessions of the 1780s and the Gadsden Purchase of 1853, the federal government acquired approximately 1.8 billion acres of land.  How to dispose of this patrimony---at what price, in what quantities, to whom, and for what purposes---became one of the central political questions of the nineteenth century, dividing parties, sections, and economic interests.

This paper traces that history in six parts.  Section~\ref{sec:framework} describes the evolving legislative framework for land sales.  Section~\ref{sec:revenue} presents data on land sales revenue and its place in federal finances.  Section~\ref{sec:parties} analyzes partisan positions on land policy.  Section~\ref{sec:railroads} examines the railroad land grants.  Section~\ref{sec:homestead} discusses the Homestead Acts and their fiscal implications.  Section~\ref{sec:conclusion} offers conclusions that connect these themes to the broader fiscal history of the early Republic.


\section{The Legislative Framework for Federal Land Sales}\label{sec:framework}

\subsection{The Land Ordinance of 1785 and the Northwest Ordinance of 1787}

The foundation of federal land policy was laid before the Constitution, under the Articles of Confederation.  The \textbf{Land Ordinance of 1785}, passed by the Congress of the Confederation on May 20, 1785, established the rectangular survey system that would eventually cover more than three-quarters of the continental United States.\footnote{The ``Point of Beginning'' for the initial survey was where Ohio, Pennsylvania, and Virginia (now West Virginia) met on the north shore of the Ohio River, near present-day East Liverpool, Ohio.}  Key provisions included:

\begin{itemize}
\item Land divided into townships of 6 miles square, each subdivided into 36 sections of 1 square mile (640 acres)
\item Section 16 of each township reserved for maintenance of public schools
\item Minimum sale price of \$1.00 per acre (raised to \$2.00 per acre in 1796)
\item Minimum purchase of 640 acres (a full section)
\item Sales at public auction with the minimum price serving as a floor
\end{itemize}

The \textbf{Northwest Ordinance of 1787} complemented the land ordinance by establishing a governance framework for the Northwest Territory (the region north of the Ohio River and east of the Mississippi).  It provided for an orderly progression from territorial status to statehood and, crucially, prohibited slavery in the territory.

These two ordinances---one providing for the survey and sale of land, the other for its governance---established the template that the federal government would apply, with modifications, across the entire western domain.

\subsection{The Harrison Land Act of 1800}

Under Treasury Secretary Alexander Hamilton and his Federalist successors, federal land policy emphasized revenue generation.  The \textbf{Land Act of 1796} raised the minimum price to \$2.00 per acre and required payment within one year, effectively limiting purchases to well-capitalized buyers.

The \textbf{Harrison Land Act of 1800}, championed by William Henry Harrison (then a delegate from the Northwest Territory), introduced important modifications:

\begin{itemize}
\item Reduced the minimum purchase to 320 acres (a half-section)
\item Permitted installment payments: one-quarter down, remainder over four years
\item Established local land offices in the western territories, reducing transaction costs for settlers
\item Maintained the \$2.00 per acre minimum
\end{itemize}

The credit provision proved double-edged.  It enabled small farmers to acquire land, but it also led to widespread overextension.  By 1820, purchasers owed the federal government approximately \$21 million in overdue installment payments---a sum exceeding annual federal revenues.

\subsection{The Land Act of 1820}

The \textbf{Land Act of 1820} (April 24, 1820) reformed the system in response to the Panic of 1819 and the accumulation of unpaid installments:

\begin{itemize}
\item \textbf{Eliminated the credit system:} Full payment required at time of purchase
\item \textbf{Reduced the minimum price:} From \$2.00 to \$1.25 per acre
\item \textbf{Reduced the minimum tract:} From 160 to 80 acres
\item \textbf{Minimum total payment:} \$100 (80 acres at \$1.25)
\end{itemize}

The \textbf{Relief Act of 1821} (February 16, 1821) further addressed the debt overhang by allowing existing purchasers to relinquish portions of their claims in exchange for forgiveness of unpaid balances.

The \$1.25 per acre price established in 1820 became the standard price for federal land for the next four decades.  It would define the arithmetic of land speculation, settlement, and revenue through the antebellum period.

\subsection{The Preemption Act of 1841}

The \textbf{Preemption Act of 1841} represented a fundamental policy shift.  It recognized the reality of ``squatting''---settlers who occupied and improved public land before it was surveyed and offered at auction---by granting them the right to purchase up to 160 acres at the minimum price of \$1.25 per acre (\$200 total) before the land was offered to other buyers at auction.

To qualify, a person had to be:
\begin{itemize}
\item At least 21 years old or head of household
\item A citizen or an immigrant who had declared intent to become a citizen
\item A resident on the land for a minimum of 14 months
\end{itemize}

The Preemption Act was significant because it subordinated the revenue motive to the settlement motive.  By guaranteeing squatters first right of purchase at the minimum price, it effectively prevented competitive bidding on improved land and reduced auction revenues.

\subsection{The Graduation Act of 1854}

The \textbf{Graduation Act of 1854}, long championed by Senator Thomas Hart Benton of Missouri, reduced the price of public lands that had remained unsold.  Lands that had been on the market for extended periods were offered at graduated (reduced) prices:

\begin{center}
\begin{tabular}{ll}
\toprule
Years Unsold & Price per Acre \\
\midrule
10--15 years & \$1.00 \\
15--20 years & \$0.75 \\
20--25 years & \$0.50 \\
25--30 years & \$0.25 \\
Over 30 years & \$0.125 \\
\bottomrule
\end{tabular}
\end{center}

The Graduation Act was repealed in 1862 when the Homestead Act made much of it moot by offering land for free.  During its brief operation, however, it significantly reduced per-acre revenues from marginal and remote lands.


\section{Land Sales Revenue and Federal Finances}\label{sec:revenue}

\subsection{Overview: Two Pillars of Federal Revenue}

Before the Civil War, the federal government relied on two principal revenue sources: customs duties (tariffs) and public land sales.  Together these accounted for virtually all federal income, as there was no federal income tax, no excise taxes after the repeal of the whiskey excise (except briefly during the War of 1812), and no significant fees or charges.

\subsection{Revenue from Land Sales: Magnitudes}

Federal land sales revenue followed a pattern driven by economic cycles, territorial acquisitions, and legislative changes:

\begin{center}
\begin{tabular}{lrrl}
\toprule
Period & Approx.\ Annual Revenue & \% of Total Revenue & Notes \\
\midrule
1790--1799 & $<$\$100,000 & $<$2\% & Minimal sales; system being established \\
1800--1810 & \$500,000--\$700,000 & 5--8\% & Harrison Act credit sales begin \\
1811--1815 & \$1--2 million & 5--10\% & War of 1812 era \\
1816--1819 & \$2--5 million & 10--20\% & Postwar land boom \\
1819--1820 & \$1.6 million & $\sim$10\% & Panic of 1819 aftermath \\
1830--1836 & \$5--25 million & 15--50\% & Jacksonian land boom \\
1836 (peak) & \$24.9 million & $\sim$48\% & Speculative frenzy \\
1837--1842 & \$2--7 million & 10--25\% & Panic of 1837 aftermath \\
1843--1852 & \$2--3 million & 5--10\% & Moderate sales \\
1853--1858 & \$5--12 million & 10--15\% & Pre-Civil War expansion \\
1862--1870 & \$1--3 million & 1--3\% & Homestead Act era \\
1880--1900 & \$4--10 million & 1--3\% & Residual sales alongside homesteads \\
\bottomrule
\end{tabular}
\end{center}

\subsection{The Spectacular Year of 1836}

The most dramatic demonstration of land sales as a revenue source came in fiscal year 1836, when land sale receipts reached \$24.9 million---nearly half of total federal revenue and more than the entire federal budget in most years.  This extraordinary sum resulted from the intersection of several forces:

\begin{enumerate}
\item \textbf{Jacksonian bank policy:} President Andrew Jackson's destruction of the Second Bank of the United States released a flood of credit from state-chartered banks
\item \textbf{Speculative fever:} Easy credit fueled massive land speculation, especially in the Mississippi Valley and the Old Northwest
\item \textbf{State investment:} States borrowed heavily to finance internal improvements (canals, roads, railroads), stimulating demand for western land
\item \textbf{Population growth:} Rapid immigration and natural increase drove westward migration
\end{enumerate}

Jackson's \textbf{Specie Circular} (July 11, 1836), which required payment for government land in gold or silver rather than banknotes, was intended to curb speculation.  Instead, it contributed to the Panic of 1837 and a sharp decline in land sales.  Revenues from land sales fell from \$24.9 million in 1836 to \$6.8 million in 1837 and just \$3.1 million in 1838.

\subsection{The Distribution Debate and the Surplus}

The surge in land revenues created a novel fiscal situation: a federal surplus so large that the national debt was entirely retired by January 1835---the only time in American history.  The surplus raised the question of what to do with excess revenues:

\begin{itemize}
\item \textbf{Henry Clay's Distribution Plan:} Whig leader Henry Clay proposed distributing federal land sale revenues to the states in proportion to their congressional representation.  This would effectively have converted federally owned western land into a national endowment fund.
\item \textbf{The Deposit Act of 1836:} Congress compromised by directing the Treasury to deposit the surplus with the states (approximately \$37 million was distributed in three installments before the Panic of 1837 halted the fourth).  This was structured as a ``deposit'' rather than a ``distribution'' to avoid constitutional objections, but the states treated the money as a gift and never returned it.
\item \textbf{Calhoun's objection:} John C. Calhoun and southern Democrats opposed distribution on the grounds that it would create incentives for high tariffs (if land revenues were given away, tariff revenues would be needed to fund the government).
\end{itemize}

\subsection{Acreages Sold}

The total disposition of the public domain is staggering in scope.  Between 1781 and the early twentieth century, the federal government disposed of approximately 1.29 billion acres through various mechanisms:

\begin{center}
\begin{tabular}{lr}
\toprule
Method of Disposition & Approximate Acres \\
\midrule
Cash sales and preemption & 300 million \\
Homestead entries & 270 million \\
Railroad land grants & 131 million \\
Military bounty warrants & 61 million \\
State grants (schools, internal improvements, swamps) & 330 million \\
Other (mineral, timber, desert entries, etc.) & $\sim$200 million \\
\midrule
\textbf{Total disposed} & $\sim$1,290 million \\
\bottomrule
\end{tabular}
\end{center}

Of this total, approximately 97\% of transfers to private ownership occurred before 1940.  Beginning in the early twentieth century, federal policy shifted from disposing of public land to retaining and managing it.

\section{Political Parties and Land Policy}\label{sec:parties}

\subsection{Federalists (1789--1800s): Revenue and Order}

The Federalist Party, led by Alexander Hamilton, viewed public lands primarily as a \textbf{source of revenue} and as a means of maintaining orderly settlement:

\begin{itemize}
\item \textbf{High minimum prices:} Federalists set the initial price at \$2.00 per acre (1796), deliberately pricing land above what many small farmers could afford
\item \textbf{Large minimum tracts:} The requirement to purchase at least 640 acres (and later 320 acres) favored well-capitalized buyers
\item \textbf{Revenue priority:} Hamilton and his allies saw land sales as a crucial complement to tariff revenues for servicing the funded national debt
\item \textbf{Orderly settlement:} Federalists opposed squatting and wanted surveyed, legally clear settlement to avoid Indian conflicts and jurisdictional disputes
\item \textbf{Eastern labor supply:} Some Federalist leaders, including Hamilton, worried that cheap western land would drain labor from eastern manufacturing, raising wages and retarding industrial development
\end{itemize}

The Federalist approach to land policy was consistent with their broader program of consolidating federal power, developing financial markets, and building a commercial republic modeled on British fiscal precedents.

\subsection{Jeffersonian Republicans (1800--1828): Yeoman Farmers}

Thomas Jefferson and his Democratic-Republican successors pursued a different vision:

\begin{itemize}
\item \textbf{Expand access:} The Harrison Land Act of 1800 (supported by Jeffersonians) reduced minimum purchases and introduced credit sales
\item \textbf{Yeoman ideal:} Jefferson's vision of an agrarian republic of independent small farmers required accessible land
\item \textbf{Louisiana Purchase (1803):} The \$15 million acquisition from France roughly doubled the national domain, vastly expanding the public lands available for sale and settlement---and was financed by borrowing from European banks, demonstrating the value of the creditworthiness Hamilton had established
\item \textbf{Revenue still important:} Despite the rhetoric of agrarian democracy, Jeffersonians continued to rely on land sales as a revenue source; Albert Gallatin, Jefferson's Treasury Secretary, used land revenues as part of his plan to retire the national debt
\item \textbf{Lower prices:} The Land Act of 1820 (under Monroe) reduced the price to \$1.25 per acre, reflecting the Jeffersonian preference for accessibility over revenue maximization
\end{itemize}

\subsection{Jacksonian Democrats (1828--1850s): Cheap Land and State Sovereignty}

Andrew Jackson and his Democratic successors pushed land policy further toward accessibility:

\begin{itemize}
\item \textbf{Preemption:} Jackson's Democrats championed preemption rights for squatters, eventually codified in the Preemption Act of 1841.  This subordinated revenue to the principle that actual settlers should have priority
\item \textbf{Graduation:} Senator Thomas Hart Benton of Missouri, Jackson's ally, pushed for graduated reduction of land prices for parcels that remained unsold, finally enacted in 1854
\item \textbf{Opposition to distribution:} Democrats opposed Henry Clay's plan to distribute land sale revenues to the states, viewing it as a scheme to justify high tariffs
\item \textbf{Specie Circular (1836):} Jackson's requirement that land be purchased with hard money aimed to curb speculation but also reflected Democratic hostility to paper money and banks
\item \textbf{State sovereignty:} Democrats generally favored ceding public lands to the states in which they were located, rather than managing them as a federal revenue source
\end{itemize}

\subsection{Whigs (1833--1856): Internal Improvements and Distribution}

The Whig Party, led by Henry Clay, Daniel Webster, and (early) Abraham Lincoln, had a distinctive land policy:

\begin{itemize}
\item \textbf{Distribution:} Clay's \textbf{American System} proposed distributing land sale revenues to the states for internal improvements (roads, canals, railroads).  The Distribution Act of 1841 briefly implemented this idea but was repealed under Tyler
\item \textbf{Maintain prices:} Whigs opposed graduation because lower land prices would reduce revenues available for distribution
\item \textbf{Internal improvements:} Whigs saw federal land as a resource to fund transportation infrastructure, not merely to be given away to settlers
\item \textbf{Tariff linkage:} Whigs explicitly connected land policy to tariff policy: if land revenues were distributed, tariffs would need to remain high to fund the government---which suited Whig protectionist interests
\item \textbf{National development:} The Whig vision emphasized using public resources (land, tariffs) for national economic development through infrastructure and education
\end{itemize}

\subsection{Republicans (1854--): Free Soil and Free Homesteads}

The Republican Party, formed in 1854, adopted the most radical position on land policy:

\begin{itemize}
\item \textbf{Free homesteads:} Republicans advocated giving public land to actual settlers for free, crystallized in the slogan ``Vote Yourself a Farm.''  The party's 1860 platform explicitly called for a homestead act
\item \textbf{Free Soil ideology:} Republicans inherited the Free Soil movement's conviction that western lands should be reserved for free white labor, not opened to slavery.  Cheap or free land would attract settlers who would establish free-state institutions
\item \textbf{Revenue sacrifice accepted:} Republicans were willing to forgo land sale revenues because they supported higher tariffs to fund the government and protect northern industry---a reversal of the Jacksonian linkage
\item \textbf{Railroad grants:} Republicans enthusiastically supported land grants to railroads as a means of developing the West and binding it to the Union
\item \textbf{Land-grant colleges:} The Morrill Act of 1862, passed alongside the Homestead Act and Pacific Railroad Act, granted federal land to states for establishing agricultural and mechanical colleges
\end{itemize}

The secession of the southern states in 1861 removed the Democratic opposition that had blocked homestead legislation.  The 37th Congress (1861--63), dominated by Republicans, passed the \textbf{Homestead Act}, the \textbf{Pacific Railroad Act}, and the \textbf{Morrill Act} in quick succession---fundamentally transforming federal land policy from revenue generation to development promotion.

\subsection{Southern Democrats and the Homestead Veto}

The antebellum politics of land policy had a crucial sectional dimension.  Southern Democrats repeatedly blocked homestead legislation because:

\begin{itemize}
\item \textbf{Slavery threat:} Free land would attract free-soil settlers to western territories, tipping the political balance against slavery
\item \textbf{Labor supply:} Cheap western land would drain poor white labor from the South, potentially raising costs for plantation owners
\item \textbf{Revenue concerns:} If land were given away free, tariff revenues would need to increase---and southern planters, as exporters, opposed high tariffs
\item \textbf{Constitutional scruples:} Some southern Democrats objected on principle to giving away national property
\end{itemize}

President James Buchanan vetoed a homestead bill in 1860.  Only after secession removed southern opposition could the Republican Congress pass the Homestead Act.


\section{Railroad Land Grants}\label{sec:railroads}

\subsection{Origins: The Land Grant Act of 1850}

The federal government began granting land to support railroad construction before the Civil War.  The \textbf{Land Grant Act of 1850} (September 20, 1850), championed by Senator Stephen A.\ Douglas of Illinois, provided 3.75 million acres of public land to the states of Illinois, Mississippi, and Alabama for railroad construction.  This act established the model:

\begin{itemize}
\item Land was granted in alternate sections (a checkerboard pattern) along the proposed railroad route
\item The railroad received alternating sections; the government retained the interspersed sections
\item The government doubled the price on its retained sections to \$2.50 per acre (from the standard \$1.25), on the theory that railroad access made adjacent land more valuable
\item Grants were made to states, which then conveyed them to railroad corporations
\end{itemize}

The principal beneficiary of the 1850 Act was the \textbf{Illinois Central Railroad}, which received approximately 2.6 million acres in Illinois and became the first major land-grant railroad.  By 1857, 21 million acres of public lands had been committed to railroad projects in the Mississippi Valley.

\subsection{The Pacific Railroad Acts (1862--1866)}

The \textbf{Pacific Railroad Act of 1862}, signed by President Lincoln on July 1, 1862, massively expanded railroad land grants and introduced a critical innovation: grants were made directly to corporations rather than to states.

\textbf{Key provisions of the 1862 Act:}

\begin{itemize}
\item Granted 10 square miles (6,400 acres) of public land for every mile of railroad constructed
\item Land was granted as 5 alternate sections per mile on each side of the track, within 10 miles of the line
\item The Central Pacific Railroad and Union Pacific Railroad were the principal charter companies
\item Additionally authorized 30-year government bonds at 6\% interest:
\begin{itemize}
\item \$16,000 per mile on the plains
\item \$32,000 per mile in the high desert between the mountain ranges
\item \$48,000 per mile through mountain ranges (limited to 300 miles at this rate)
\end{itemize}
\item Government bonds were secured by a lien on the railroads and their fixtures
\end{itemize}

\textbf{The 1864 Amendment} doubled the land grant to 20 sections (12,800 acres) per mile and subordinated the government's bond lien to the railroad companies' own first-mortgage bonds, improving the companies' ability to raise private capital.

\subsection{Total Railroad Land Grants: Magnitude and Value}

From 1850 to 1871 (when Congress ended the practice of new railroad land grants), the federal government granted more than \textbf{175 million acres} of public land to railroad companies---an area larger than the state of Texas and more than one-tenth of the entire United States.

Including state grants made from federally donated land, the total approaches \textbf{131 million acres} that were ultimately patented (with some grants being forfeited for failure to build the required railroad).

\begin{center}
\begin{tabular}{lrr}
\toprule
Railroad & Approximate Acres Granted & Federal Bonds Issued \\
\midrule
Union Pacific & 12,000,000 & \$27,236,512 \\
Central Pacific & 8,000,000 & \$27,855,680 \\
Northern Pacific & 39,000,000 & None \\
Southern Pacific (via Texas Pacific) & 18,000,000 & None \\
Atchison, Topeka \& Santa Fe & 17,000,000 & None \\
Illinois Central & 2,600,000 & None \\
Other railroads & $\sim$35,000,000 & Varies \\
\bottomrule
\end{tabular}
\end{center}

All government bond loans were eventually repaid in full with interest by the railroad companies.

\subsection{Railroad Land Sales: Prices and Values}

The railroads sold their granted lands to settlers and investors at prices that varied widely depending on location, soil quality, and proximity to the track:

\begin{center}
\begin{tabular}{ll}
\toprule
Category & Typical Price per Acre \\
\midrule
Prime agricultural land near stations & \$5.00--\$15.00 \\
Good farmland within 5 miles of track & \$3.00--\$8.00 \\
Average land 5--10 miles from track & \$2.00--\$5.00 \\
Marginal or remote land & \$1.25--\$3.00 \\
Arid or desert land & \$0.50--\$1.25 \\
\bottomrule
\end{tabular}
\end{center}

The \textbf{Illinois Central Railroad} provides the best-documented case.  Between 1855 and 1890, the IC sold its 2.6 million acres of granted land for a total of approximately \$16.5 million, an average of roughly \$6.30 per acre---more than five times the government minimum price of \$1.25.  The railroad actively promoted settlement through advertising campaigns in the eastern United States and in Europe, particularly in Scandinavia, Germany, and Britain.

The \textbf{Northern Pacific Railroad}, which received the largest grant (39 million acres), sold its lands at lower average prices---approximately \$2.50--\$4.00 per acre---reflecting the more remote and arid character of much of its grant in the northern Great Plains and Pacific Northwest.

\subsection{The Checkerboard Logic and Revenue Implications}

The alternate-sections grant pattern had a specific revenue logic:

\begin{enumerate}
\item The government gave away every other section along the railroad right-of-way
\item But it doubled the price on its retained sections from \$1.25 to \$2.50 per acre
\item In theory, the government recovered equivalent revenue from selling half the land at double the price
\item In practice, the calculation was more favorable because railroad access genuinely increased land values, so the retained sections were often worth more than \$2.50
\end{enumerate}

Whether the railroad land grants represented a net subsidy to the railroads or a net benefit to the government has been extensively debated.  The railroads provided reduced rates for government freight and mail (a condition of the grants) that partially offset the value of the land.  A 1944 study estimated that these government rate reductions saved the federal government approximately \$500 million between 1850 and 1945.

\subsection{Political Economy of Railroad Grants}

Railroad land grants involved all the classic features of political economy:

\begin{itemize}
\item \textbf{Lobbying and corruption:} Railroad promoters aggressively lobbied Congress.  Cr\'edit Mobilier of America, the construction company behind the Union Pacific, distributed stock to members of Congress in what became one of the great scandals of the Gilded Age (exposed in 1872)
\item \textbf{Sectional conflict:} The route of the transcontinental railroad was debated for a decade before the Civil War; southern Democrats blocked a central or northern route, and only secession made a central route politically feasible
\item \textbf{End of grants (1871):} Growing public opposition to railroad land grants, fueled by perceptions of corruption and monopoly power, led Congress to end new grants in 1871, though existing grants continued to be executed
\item \textbf{Forfeiture movement (1880s--1890s):} Reformers demanded that unearned railroad land grants (where the railroad had not built the required line) be forfeited back to the government.  Approximately 35 million acres of grants were forfeited between 1890 and 1940
\end{itemize}


\section{The Homestead Acts and Revenue Consequences}\label{sec:homestead}

\subsection{The Homestead Act of 1862}

The \textbf{Homestead Act of 1862} (May 20, 1862), signed by President Lincoln, represented the most fundamental shift in federal land policy since the Land Ordinance of 1785.  It replaced the principle of \textbf{revenue maximization} with the principle of \textbf{free distribution to actual settlers}:

\textbf{Key provisions:}
\begin{itemize}
\item Any adult citizen (or declared intent to become citizen) who had never borne arms against the United States could claim 160 acres of surveyed public land
\item The claimant had to reside on the land and cultivate it for five years
\item After five years of continuous residence and improvement, the claimant received a patent (deed) for a nominal fee (\$10 filing fee plus \$8 final fee)
\item Alternatively, after six months of residence, the claimant could ``commute'' the homestead entry by purchasing the land at \$1.25 per acre
\item The land was exempt from seizure for debts incurred before the patent was issued
\item Women were eligible, as were immigrants who had filed a declaration of intent to become citizens
\end{itemize}

\subsection{Revenue Impact}

The fiscal consequences of the Homestead Act were dramatic.  Land sales revenue, which had been a significant component of federal income, declined sharply:

\begin{center}
\begin{tabular}{lrl}
\toprule
Period & Land Revenue (annual avg.) & Context \\
\midrule
1850--1860 & \$3--8 million & Pre-Homestead Act \\
1862--1870 & \$1--3 million & Homesteading begins \\
1870--1880 & \$2--4 million & Mixed sales and homesteads \\
1880--1890 & \$5--10 million & Some revival in cash sales \\
1890--1900 & \$2--4 million & Frontier closing \\
\bottomrule
\end{tabular}
\end{center}

However, the decline in land revenues was more than offset by three factors:

\begin{enumerate}
\item \textbf{Higher tariffs:} The Republican Congresses of the Civil War era raised tariffs dramatically (the Morrill Tariff of 1861 and subsequent increases), replacing land revenues in the federal budget
\item \textbf{Income tax and excises:} The Revenue Act of 1861 introduced the first federal income tax; internal revenue (excises on whiskey, tobacco, manufactures, and income taxes) became a major revenue source during and after the Civil War
\item \textbf{Economic growth:} Settlement of homesteaded land expanded the tax base through increased trade, manufacturing, and consumption of imported (tariffed) goods
\end{enumerate}

\subsection{Homesteading by the Numbers}

Between 1862 and 1934, the Homestead Act and its successors distributed land on an enormous scale:

\begin{center}
\begin{tabular}{lr}
\toprule
Statistic & Value \\
\midrule
Total homestead entries filed & 1.6 million \\
Total acres distributed & 270 million \\
Percentage of total U.S.\ land area & $\sim$10\% \\
Completion rate (``proving up'') & $\sim$40\% \\
Peak year for entries (acres) & 1910 (18.3 million acres) \\
Descendants of homesteaders (as of 2017) & $\sim$93 million Americans \\
\bottomrule
\end{tabular}
\end{center}

\subsection{Related Legislation}

The Homestead Act was supplemented by several related measures:

\begin{itemize}
\item \textbf{Timber Culture Act of 1873:} Granted an additional 160 acres to homesteaders who planted at least 10 acres of trees, allowing a total homestead of 320 acres
\item \textbf{Desert Land Act of 1877:} Allowed purchase of 640 acres of arid land at \$1.25 per acre, provided the claimant irrigated it within three years
\item \textbf{Kinkaid Act of 1904:} Expanded homesteads in western Nebraska to 640 acres, recognizing that semi-arid land required larger holdings
\item \textbf{Enlarged Homestead Act of 1909:} Increased homesteads to 320 acres for dryland farming areas in the Great Plains
\item \textbf{Stock-Raising Homestead Act of 1916:} Allowed 640-acre homesteads for ranching purposes in arid regions
\end{itemize}

Each expansion reflected the reality that 160 acres---adequate for eastern agriculture---was insufficient in the semi-arid and arid West.

\subsection{Commutation and Cash Sales}

Despite the Homestead Act's offer of free land, cash sales continued.  Many homesteaders used the ``commutation'' provision to purchase their claims after six months by paying \$1.25 per acre, because:

\begin{itemize}
\item It was faster than waiting five years
\item It gave them clear title to sell, mortgage, or otherwise use as collateral
\item Speculators sometimes used dummy entrymen to file homestead claims and then commute them
\end{itemize}

Additionally, certain categories of land---mineral lands, timber lands, and some railroad forfeiture lands---were still available only by cash purchase.  These residual cash sales provided some continued revenue to the Treasury even after the Homestead Act.

\subsection{The End of the Frontier and the Shift to Retention}

Frederick Jackson Turner famously declared the frontier ``closed'' in 1890, based on the Census Bureau's determination that there was no longer a continuous line of unsettled territory.  In practice, homesteading continued well into the twentieth century, peaking in 1910 and declining sharply after 1935.

The \textbf{Federal Land Policy and Management Act of 1976} formally ended homesteading in the lower 48 states (Alaska continued until 1986).  The last homestead patent was issued to Ken Deardorff for 80 acres on the Stony River in southwestern Alaska; he fulfilled the requirements in 1979 but did not receive his deed until May 1988.

By this time, federal policy had shifted decisively from disposal to retention.  Today, the federal government retains approximately 640 million acres---28\% of the total U.S.\ land area---managed by the Bureau of Land Management, Forest Service, Fish and Wildlife Service, and National Park Service.


\section{Synthesis: Land Policy in the Long View of Fiscal History}\label{sec:conclusion}

\subsection{From Revenue Source to Development Tool}

The arc of federal land policy from 1785 to 1900 traces a fundamental transformation:

\begin{enumerate}
\item \textbf{Revenue era (1785--1840s):} Land sales were a primary revenue source, second only to tariffs.  Policy emphasized high prices, minimum tract sizes, and orderly sales to maximize receipts
\item \textbf{Transition era (1841--1862):} Preemption, graduation, and the growing political power of western settlers progressively eroded the revenue model in favor of accessible settlement
\item \textbf{Development era (1862--1900):} The Homestead Act, railroad grants, and land-grant colleges shifted the purpose of the public domain from generating revenue to promoting economic development through settlement, transportation infrastructure, and education
\end{enumerate}

\subsection{Fiscal Substitution}

The change in land policy required a parallel transformation of the federal revenue system:

\begin{center}
\begin{tabular}{lll}
\toprule
Revenue Source & Pre-Civil War & Post-Civil War \\
\midrule
Tariffs & 80--90\% of revenue & 40--50\% of revenue \\
Land sales & 5--20\% of revenue & 1--5\% of revenue \\
Internal revenue (excises, income tax) & 0--5\% & 30--50\% \\
\bottomrule
\end{tabular}
\end{center}

The willingness to sacrifice land revenues was politically feasible precisely because the Civil War had created new revenue instruments---particularly the income tax and internal excises---that provided abundant alternative sources.  The Republicans who supported free homesteads were the same Republicans who supported high tariffs and (temporarily) income taxes.

\subsection{Continuity with Hamilton's Project}

Ironically, the massive giveaway of public land in the 1860s--1900s was possible only because Hamilton's earlier fiscal revolution had succeeded.  Hamilton's funding system:

\begin{itemize}
\item Established federal creditworthiness, enabling the government to borrow for territorial acquisitions (Louisiana Purchase, Mexican Cession) that created the public domain
\item Demonstrated that tariff revenues could reliably service the national debt, reducing dependence on land revenues
\item Created financial markets and institutions that facilitated the complex railroad land grant and bond systems
\item Established the precedent of federal fiscal supremacy over the states, which made federal (rather than state) disposal of western lands politically sustainable
\end{itemize}

The story of federal land sales is thus a continuation of the fiscal state-building narrative that begins with Hamilton's confrontation of the land speculators and public creditors in 1790.

\subsection{Magnitudes in Perspective}

To appreciate the scale of the public domain's disposition:

\begin{itemize}
\item The total public domain at its maximum extent was approximately \textbf{1.84 billion acres}
\item Of this, approximately \textbf{1.29 billion acres} were transferred out of federal ownership
\item At the standard price of \$1.25 per acre, the full public domain would have been worth approximately \textbf{\$2.3 billion} in nominal terms---more than 25 times the peak national debt under Hamilton's funding system
\item Much of this value was given away (homesteads, railroad grants, state grants, military bounties) rather than sold
\item The total cumulative cash receipts from land sales between 1785 and 1880 were approximately \textbf{\$200 million}---roughly 10\% of the theoretical value at the minimum price
\end{itemize}

The public domain was, in effect, the largest asset on the federal balance sheet for most of the nineteenth century, and its disposition---whether for revenue, settlement, transportation, or education---was the single most consequential set of fiscal decisions the federal government made between Hamilton's funding system and the introduction of the income tax.

\subsection{Lessons for Fiscal History}

The land sales episode illustrates several principles relevant to the broader fiscal history:

\begin{enumerate}
\item \textbf{Revenue sources shape political coalitions:} Dependence on tariffs and land sales created different political dynamics than dependence on income taxes.  Revenue sources are never politically neutral.

\item \textbf{Asset disposition involves commitment problems:} Once the government began giving land away (preemption, homesteads), it was politically impossible to return to revenue-maximizing sales.  Each concession created constituencies demanding further concessions.

\item \textbf{Infrastructure subsidies have complex fiscal effects:} Railroad land grants sacrificed direct land revenues but generated indirect fiscal benefits through economic development, increased trade (tariff revenues), and enhanced government bond values.

\item \textbf{Distributional conflicts drive policy:} The battle over land policy was ultimately a battle among eastern manufacturers (who wanted labor to stay east), western settlers (who wanted free or cheap land), southern planters (who wanted land open to slavery), and speculators (who wanted to profit from intermediation).

\item \textbf{Fiscal capacity enables policy flexibility:} The ability to sacrifice land revenues for development purposes depended on having alternative revenue sources.  Hamilton's tariff and debt system made later land policy liberalization financially feasible.
\end{enumerate}

\bibliographystyle{plainnat}
\bibliography{fiscal_history}

\end{document}
